\usepackage{amsmath,amssymb,amsthm}
\usepackage{cases} 
\usepackage{tikz-cd}        %diagram

\usepackage[
backend=biber,
style=numeric,%alphabetic,
sorting=ynt
]{biblatex}
\addbibresource{mybibliography.bib}


\usepackage{hyperref}
\hypersetup{colorlinks}
\hypersetup{
    bookmarksopen=true,     % show bookmarks bar?
    unicode=false,          % non-Latin characters in Acrobat’s bookmarks
    pdftoolbar=true,        % show Acrobat’s toolbar?
    pdfmenubar=true,        % show Acrobat’s menu?
    pdffitwindow=false,     % window fit to page when opened
    pdfstartview={FitH},    % fits the width of the page to the window
    pdftitle={My title},    % title
    pdfauthor={Author},     % author
    pdfsubject={Subject},   % subject of the document
    pdfcreator={Creator},   % creator of the document
    pdfproducer={Producer}, % producer of the document
    pdfkeywords={keyword1} {key2} {key3}, % list of keywords
    pdfnewwindow=true,      % links in new PDF window
    colorlinks=true,       % false: boxed links; true: colored links
    linkcolor=red,          % color of internal links (change box color with linkbordercolor)
    citecolor=magenta,        % color of links to bibliography
    filecolor=violet,      % color of file links
    urlcolor=blue      % color of external links
}
%%%%%%%%%%%%%
\textwidth=14.5cm
\textheight=21.5cm
\oddsidemargin=0.8cm
\evensidemargin=0.8cm
\topmargin=1cm
%%%%%%%%%%%%%
\newtheorem{theorem}{Theorem}[section]
\newtheorem{lemma}[theorem]{Lemma}
\newtheorem{corollary}{Corollary}[theorem]
\newtheorem{definition}[theorem]{Definition}
\newtheorem{proposition}[theorem]{Proposition}
\newtheorem{example}[theorem]{Example}
\newtheorem{question}[theorem]{Question}
\newtheorem*{remark}{Remark}
%%%%%%%%%%%%%
